\chapter{BUN (Blood Urea Nitrogen)}

\section*{What Is This Test?}
The blood urea nitrogen (BUN) test measures the amount of urea nitrogen in the blood. Urea is the primary nitrogenous end product of protein and amino acid catabolism, produced in the liver through the urea cycle (ornithine cycle). Dietary protein and endogenous protein breakdown release amino acids, which are deaminated to form ammonia. Ammonia is toxic to the central nervous system and must be rapidly converted to non-toxic urea in the liver via the urea cycle involving five enzymes: carbamoyl phosphate synthetase I, ornithine transcarbamylase, argininosuccinate synthetase, argininosuccinate lyase, and arginase. Urea is then released into the bloodstream and almost exclusively excreted by the kidneys (approximately 85--90\%), with the remainder excreted through the GI tract and skin. BUN concentration reflects the balance between hepatic urea production and renal urea excretion. It is filtered freely at the glomerulus, and approximately 40--50\% is passively reabsorbed in the proximal tubule along with water. In states of volume depletion, increased water and urea reabsorption from increased ADH activity causes a disproportionate rise in BUN relative to creatinine (BUN:Cr ratio >20:1), which is the hallmark of prerenal azotemia. BUN is measured as part of the basic metabolic panel (BMP) and is always interpreted in conjunction with serum creatinine.

\section*{Alternative Names}
Blood Urea Nitrogen; Urea Nitrogen; Serum Urea; BUN; Urea; Blood Urea; Plasma Urea Nitrogen

\section*{Principle}
BUN is measured on automated chemistry analyzers using enzymatic methods. The most common method uses urease, which hydrolyzes urea to ammonia and carbon dioxide: Urea + H\textsubscript{2}O $\rightarrow$ 2 NH\textsubscript{3} + CO\textsubscript{2}. The ammonia produced is then quantified by one of several coupled enzymatic reactions: (1) the glutamate dehydrogenase (GLDH) method: ammonia + alpha-ketoglutarate + NADH $\rightarrow$ glutamate + NAD\textsuperscript{+}; the decrease in absorbance at 340 nm is proportional to urea concentration; or (2) the Berthelot reaction with phenol and hypochlorite producing a blue indophenol chromophore. The urease-GLDH method is the most widely used in modern analyzers (Roche Cobas, Abbott Architect, Siemens Atellica). Some point-of-care devices use urease-based potentiometric methods with an ammonium-selective electrode. Results are reported as BUN in mg/dL in the United States, which measures the nitrogen portion of the urea molecule only. In many other countries, results are reported as urea in mmol/L (BUN in mg/dL $\times$ 2.14 = urea in mg/dL; urea in mg/dL / 2.8 = BUN in mmol/L). This test is performed on serum separator tubes using enzymatic spectrophotometric methods, not by hematology analyzers.

\section*{History}
The urea cycle was elucidated by Hans Krebs and Kurt Henseleit in 1932 (the first metabolic cycle to be discovered, predating the Krebs/TCA cycle by 5 years). The urease method for urea measurement was developed in the early 20th century. The Berthelot reaction for ammonia detection dates to 1859. The automated methods using urease-GLDH coupled reactions were introduced in the 1960s-1970s and remain the standard. The clinical utility of the BUN:creatinine ratio for differentiating prerenal from intrinsic renal disease was established by Dossetor in the 1960s \cite{dossetor1966creatininemia} and refined by subsequent work \cite{morgan1977plasma}.

\section*{Why This Test Is Ordered}
\begin{itemize}
  \item Evaluation of renal function as part of the basic metabolic panel (BMP): BUN is measured simultaneously with creatinine, and the BUN:Cr ratio is essential for classifying the etiology of azotemia
  \item Assessment of prerenal azotemia in dehydrated patients: BUN:Cr ratio >20:1 suggests prerenal cause (dehydration, heart failure, shock, GI bleeding, high protein intake, corticosteroid use, tetracycline)
  \item Evaluation of gastrointestinal bleeding: blood in the GI tract is digested by gut bacteria, and the resulting amino acids are absorbed and converted to urea, causing a disproportionate rise in BUN (BUN:Cr >30:1, even >36:1) with normal creatinine. This is a classic clinical clue for upper GI bleeding
  \item Monitoring of renal function during nephrotoxic drug therapy: aminoglycosides, cisplatin, amphotericin B, NSAIDs, vancomycin, contrast dye, ACE inhibitors/ARBs
  \item Assessment of nutritional status and protein intake in hospitalized patients, burn patients, and those on TPN
  \item Evaluation of hepatic synthetic function: severely low BUN (<5 mg/dL) in a patient with liver disease suggests impaired hepatic urea synthesis and poor hepatic function
  \item Monitoring of dialysis adequacy: BUN reduction ratio (URR) = (pre-dialysis BUN -- post-dialysis BUN) / pre-dialysis BUN $\times$ 100; target URR >65\%
  \item Assessment of catabolic states: trauma, sepsis, burns, surgery, corticosteroid therapy, and tetracycline therapy increase protein catabolism and BUN
\end{itemize}

\section*{Primary vs Secondary Test}
BUN is a primary screening test as an integral component of the basic metabolic panel (BMP). It is essential for the initial approach to any patient with suspected renal dysfunction and for monitoring hospitalized patients. Interpretation always requires simultaneous serum creatinine and clinical volume status assessment.

\section*{How This Test Is Performed}
A healthcare professional draws blood from a vein into a serum separator tube (gold-top) or lithium heparin tube (green-top). Approximately 3--5 mL of blood is collected. The sample is centrifuged and analyzed on an automated chemistry analyzer using the urease-GLDH enzymatic method. Results are typically available within 30--60 minutes. No special handling is required beyond standard laboratory practice.

\section*{What Preparation Is Needed}
No fasting is required for BUN alone. However, several factors affect BUN and should be considered: (1) High protein intake (meat-heavy meal) within 24 hours can significantly elevate BUN. A diet with >2 g/kg/day protein will increase BUN regardless of renal function. (2) Certain medications increase BUN independent of renal function: corticosteroids (catabolic effect), tetracyclines (anti-anabolic effect), and nephrotoxic drugs (actually impair GFR). (3) Strenuous exercise may cause a slight rise from protein catabolism. (4) Advanced age is associated with declining GFR, so BUN rises with age (approximately 1 mg/dL per decade after age 40 even in healthy individuals). (5) Inform the provider of all medications, especially diuretics, NSAIDs, ACE inhibitors, and recent contrast dye exposure.

\section*{Contraindications and Precautions}
\begin{itemize}
  \item No absolute contraindications. Important pre-analytical and clinical considerations:
  \item (1) BUN is affected by numerous non-renal factors and must always be interpreted with serum creatinine.
  \item (2) The BUN:creatinine ratio is useful but has limitations: it assumes stable creatinine and BUN production rates, which may not hold in critical illness.
  \item (3) Ammonium-containing anticoagulants (ammonium heparin) cannot be used because they artificially elevate BUN from ammonia contamination of the urease method.
  \item (4) Sodium fluoride/potassium oxalate tubes (gray-top) can inhibit urease activity and produce falsely low BUN in some assays.
  \item (5) Hemolysis does not significantly affect BUN measurement.
  \item (6) Drugs interfering with the urease method: fluorides, certain antibiotics at very high concentrations, but these are typically not clinically relevant at standard doses.
  \item (7) Infants have lower BUN (3--12 mg/dL) due to higher anabolic state and lower protein intake relative to body weight.
  \item (8) Pregnancy: BUN decreases due to increased GFR (glomerular hyperfiltration), typically 5--10 mg/dL in the second and third trimesters.
\end{itemize}
\section*{Risks}
Minimal risk from standard venipuncture. Mild pain, bruising, or hematoma at the puncture site resolves within 1--2 days. Rare complications: infection, nerve injury, vasovagal syncope. Blood volume drawn (3--5 mL) is negligible.

\section*{What Happens After the Test}
Results available within 30--60 minutes. The BUN is always interpreted with creatinine and the BUN:Cr ratio: Normal ratio 10:1 to 20:1. Ratio >20:1: prerenal azotemia (dehydration, heart failure, cirrhosis, shock, GI bleeding, high protein intake, corticosteroid use, tetracycline, burns, severe infection). Ratio 10:1 to 15:1: intrinsic renal disease (acute tubular necrosis, glomerulonephritis, interstitial nephritis). Ratio <10:1: acute tubular necrosis with low BUN production (hepatic insufficiency), malnutrition, pregnancy, SIADH (dilutional), overaggressive IV fluids. Ratio >20:1 with BUN >100 mg/dL and normal creatinine strongly suggests GI bleeding. Serial BUN monitoring is essential for assessing response to fluid resuscitation in prerenal states, monitoring nephrotoxicity, and evaluating dialysis adequacy. BUN >100 mg/dL generally indicates significant renal impairment or severe prerenal state requiring intervention.

\section*{Understanding Results}

\subsection*{Below Normal}
\begin{itemize}
  \item \textbf{Low BUN (<6 mg/dL):} Relatively uncommon; often iatrogenic or reflecting decreased hepatic urea production.
  \item \textbf{Severe liver disease:} Cirrhosis, fulminant hepatic failure, Reye syndrome, urea cycle disorders (inborn errors). The liver cannot produce urea, so BUN remains low regardless of renal function. In decompensated cirrhosis, BUN may be 2--5 mg/dL despite normal or impaired renal function.
  \item \textbf{Pregnancy:} Physiological expansion of plasma volume and increased GFR cause dilution and increased renal urea clearance. BUN typically 5--10 mg/dL in the second and third trimesters.
  \item \textbf{SIADH:} Water retention causes dilution of all solutes including urea.
  \item \textbf{Malnutrition and protein restriction:} Severely decreased dietary protein reduces urea production. Seen in anorexia nervosa, kwashiorkor, cachexia, and prolonged illness.
  \item \textbf{Anabolic states:} Growth hormone therapy, anabolic steroid use, recovery from major illness---increased protein synthesis reduces urea generation.
  \item \textbf{Overaggressive IV fluid administration:} Dilutional effect plus increased renal urea clearance from volume expansion.
  \item \textbf{Nephrotic syndrome:} Loss of urea (a small molecule) through the damaged glomerular basement membrane into urine.
\end{itemize}

\subsection*{Above Normal}
\begin{itemize}
  \item \textbf{Prerenal azotemia (BUN:Cr >20:1):} Most common cause of elevated BUN. Reduced effective circulating volume decreases renal perfusion and GFR. The kidney responds by increasing water and urea reabsorption (ADH and aldosterone-mediated). Causes: volume depletion (vomiting, diarrhea, diuretics, hemorrhage, burns, third-space losses), decreased cardiac output (CHF, cardiogenic shock), systemic vasodilation (sepsis, anaphylaxis, cirrhosis), renal artery stenosis (bilateral), drugs reducing glomerular perfusion (NSAIDs, ACE-I/ARBs in renal artery stenosis), hepatorenal syndrome.
  \item \textbf{Gastrointestinal bleeding (BUN:Cr >30:1, often >36:1):} Blood proteins in the GI lumen are digested to amino acids, absorbed, and converted to urea by the liver. This produces a disproportionate rise in BUN with normal or near-normal creatinine. An upper GI bleed is a classic cause of elevated BUN without creatinine elevation. Absence of this finding does not rule out GI bleeding.
  \item \textbf{High protein intake:} Enteral or parenteral nutrition with high protein content; high-protein diet (>2 g/kg/day); massive tissue necrosis (burns, trauma, rhabdomyolysis); GI hemorrhage. Corticosteroid therapy and tetracyclines increase protein catabolism, elevating BUN production.
  \item \textbf{Intrinsic renal disease (BUN:Cr usually 10:1 to 15:1):} Acute tubular necrosis (ATN): ischemic (prolonged prerenal state, shock) or nephrotoxic (aminoglycosides, cisplatin, amphotericin B, contrast dye, myoglobin, hemoglobin). Glomerulonephritis (post-streptococcal, IgA, lupus nephritis, anti-GBM disease, ANCA vasculitis). Acute interstitial nephritis (drugs: NSAIDs, penicillins, cephalosporins, sulfonamides, PPIs, rifampin; infections; autoimmune). Chronic kidney disease (diabetic nephropathy, hypertensive nephrosclerosis, polycystic kidney disease).
  \item \textbf{Postrenal azotemia (BUN rises proportionally with creatinine):} Obstruction to urine flow: benign prostatic hyperplasia, prostate cancer, bladder outlet obstruction, bilateral ureteral obstruction (stones, tumors, retroperitoneal fibrosis, pelvic malignancy), neurogenic bladder. Postrenal azotemia initially has a BUN:Cr ratio like prerenal (>20:1) due to increased tubular pressure causing back-diffusion of urea. With prolonged obstruction, ratio falls toward 10:1 as tubule damage progresses.
  \item \textbf{Critical values:} BUN >100 mg/dL generally warrants nephrology consultation and consideration of dialysis. BUN >200--300 mg/dL is often seen in patients requiring emergent dialysis, particularly when accompanied by uremic symptoms (pericarditis, encephalopathy, bleeding, nausea/vomiting).
\end{itemize}

\section*{Normal Results}
\begin{itemize}
  \item \textbf{Adults:} 6--20 mg/dL (2.1--7.1 mmol/L as urea)
  \item \textbf{Adults over 60 years:} 8--23 mg/dL (gradual increase with age due to declining GFR)
  \item \textbf{Children (1--16 years):} 7--20 mg/dL
  \item \textbf{Infants (1--12 months):} 4--19 mg/dL
  \item \textbf{Newborns (0--7 days):} 3--12 mg/dL
  \item \textbf{Pregnancy (2nd and 3rd trimester):} 5--10 mg/dL
  \item \textbf{BUN:Creatinine ratio:} 10:1 to 20:1
  \item \textbf{Conversion:} BUN (mg/dL) ÷ 2.8 = BUN (mmol/L); BUN (mg/dL) $\times$ 2.14 = Urea (mg/dL)
\end{itemize}

\section*{Follow-up Tests}
\begin{itemize}
  \item \textbf{Serum creatinine with eGFR:} Essential for comprehensive renal function assessment and CKD staging. Must be ordered simultaneously with BUN.
  \item \textbf{BUN:Creatinine ratio:} Calculated automatically. >20:1 suggests prerenal, GI bleed, catabolic state. 10:1--15:1 suggests intrinsic renal. <10:1 suggests hepatic insufficiency, malnutrition, pregnancy.
  \item \textbf{Urinalysis with microscopy:} Presence of muddy brown casts, renal tubular epithelial cells strongly suggests ATN. Dysmorphic RBCs and RBC casts suggest glomerulonephritis. WBC casts suggest interstitial nephritis. Eosinophiluria suggests drug-induced interstitial nephritis.
  \item \textbf{Urine sodium (UNa) and fractional excretion of sodium (FENa):} FENa <1\% suggests prerenal azotemia; FENa >2--3\% suggests ATN. Caution: FENa is unreliable in patients on diuretics (use FEUrea instead). FEUrea <35\% suggests prerenal; >50\% suggests ATN.
  \item \textbf{Serum electrolytes (BMP/CMP):} Hyperkalemia and metabolic acidosis are common in renal failure.
  \item \textbf{Renal ultrasound:} Assess kidney size, cortical thickness, echogenicity, and hydronephrosis. Small, shrunken kidneys (<8--9 cm) suggest chronic irreversible disease. Hydronephrosis indicates postrenal obstruction.
  \item \textbf{Serum albumin, INR, and bilirubin:} Assess hepatic synthetic function in low-BUN states with concomitant liver disease.
  \item \textbf{Upper endoscopy:} If GI bleeding suspected based on BUN:Cr ratio.
  \item \textbf{Creatinine kinase (CK):} If rhabdomyolysis is the suspected cause of ATN.
  \item \textbf{24-hour urine collection for creatinine clearance:} More accurate GFR assessment than eGFR in patients with extremes of muscle mass, amputees, or rapidly changing renal function.
\end{itemize}

\section*{References}
\bibliographystyle{unsrtnat}
\bibliography{references}

\clearpage
\section*{Regulatory Status and Authorization Date}
This chapter describes a test category, not a specific commercial product. FDA, CDSCO, EU IVDR, and MHRA approval or registration dates therefore cannot be assigned without the assay manufacturer, product name, instrument, and intended use. For laboratory-developed tests, an individual agency approval date may not exist. See the Preface for the interpretation of regulatory status.

\clearpage
